\section{Phương trình và chuỗi phản ứng của kim loại}

\subsection{Luyện tập viết phương trình hóa học}

Trước khi làm chuỗi phản ứng, cần viết chắc các phản ứng đơn lẻ. Với mỗi phản ứng, hãy xác định loại chất ban đầu, dự đoán sản phẩm rồi cân bằng phương trình.

\begin{tasks}(2)
\task $\ce{Mg + O2 ->}$
\task $\ce{Al + O2 ->}$
\task $\ce{Fe + Cl2 ->}$
\task $\ce{Fe + S ->}$
\task $\ce{Zn + HCl ->}$
\task $\ce{Fe + H2SO4 ->}$
\task $\ce{Fe + CuSO4 ->}$
\task $\ce{Cu + AgNO3 ->}$
\task $\ce{FeCl2 + NaOH ->}$
\task $\ce{FeCl3 + NaOH ->}$
\task $\ce{Fe(OH)2 + O2 + H2O ->}$
\task $\ce{Fe(OH)3 ->[t^\circ]}$
\end{tasks}

\begin{ghinho}{Cách làm chuỗi phản ứng}
Với mỗi mũi tên trong chuỗi:
\begin{enumerate}
\item xác định chất đầu và chất cuối;
\item chọn thuốc thử hoặc điều kiện phù hợp;
\item viết phương trình hóa học;
\item kiểm tra sản phẩm tạo thành có đúng chất cần tìm hay không.
\end{enumerate}
Không giải cả chuỗi cùng lúc; hãy xử lí từng mũi tên.
\end{ghinho}

\begin{vidu}
Hoàn thành chuỗi phản ứng:
\[
\ce{Fe -> FeCl2 -> Fe(OH)2 -> Fe(OH)3 -> Fe2O3}.
\]
\end{vidu}
\begin{loigiai}
\[
\ce{Fe + 2HCl -> FeCl2 + H2 ^},
\]
\[
\ce{FeCl2 + 2NaOH -> Fe(OH)2 v + 2NaCl},
\]
\[
\ce{4Fe(OH)2 + O2 + 2H2O -> 4Fe(OH)3},
\]
\[
\ce{2Fe(OH)3 ->[t^\circ] Fe2O3 + 3H2O}.
\]

\textbf{Vậy,} chuỗi được thực hiện bằng bốn phản ứng trên.
\end{loigiai}

\begin{tuluyen}
Hoàn thành các chuỗi phản ứng sau:
\begin{tasks}(1)
\task $\ce{Cu -> CuO -> CuCl2 -> Cu(OH)2 -> CuO}$.
\task $\ce{Fe -> FeCl3 -> Fe(OH)3 -> Fe2O3 -> Fe}$.
\task $\ce{Zn -> ZnO -> ZnCl2 -> Zn(OH)2 -> ZnO}$.
\end{tasks}
\dapso{Viết các PTHH tương ứng; ở bước \ce{Fe2O3 -> Fe} có thể dùng CO ở nhiệt độ cao.}
\end{tuluyen}

\subsection{Xác định chất từ thành phần phần trăm}

Một số đề cho thành phần phần trăm của nguyên tố trong oxide hoặc muối rồi yêu cầu xác định công thức. Cách làm:
\begin{enumerate}
\item đặt công thức tổng quát phù hợp;
\item lập biểu thức phần trăm khối lượng;
\item tìm chỉ số nguyên phù hợp;
\item dùng chất vừa xác định để hoàn thành chuỗi phản ứng.
\end{enumerate}

\begin{vidu}
Một oxide sắt chứa $72.4\%$ Fe theo khối lượng. Xác định oxide, biết oxide thuộc một trong ba chất \ce{FeO}, \ce{Fe2O3}, \ce{Fe3O4}.
\end{vidu}
\begin{loigiai}
Tính phần trăm Fe trong từng oxide:
\[
\%Fe_{\ce{FeO}}=\frac{56}{72}\times100\%\approx77.8\%,
\]
\[
\%Fe_{\ce{Fe2O3}}=\frac{112}{160}\times100\%=70.0\%,
\]
\[
\%Fe_{\ce{Fe3O4}}=\frac{168}{232}\times100\%\approx72.4\%.
\]

\textbf{Vậy,} oxide cần tìm là \ce{Fe3O4}.
\end{loigiai}
